\documentclass[aspectratio=169,11pt]{beamer}

% Shared style, ported from the UF Animal Sciences 2026 talk (ib-tutorials.sty at the repo root).
\makeatletter\def\input@path{{./}{../}}\makeatother
\usepackage{ib-tutorials}

\usepackage[utf8]{inputenc}
\usepackage{graphicx}
\usepackage{amsmath}

\title[IBM backup]{Immersed Boundary Method Backup Slides for Exams and Talks}
\author{Laura A. Miller}
\institute{Department of Mathematics, University of Arizona}
\date{Fluid Dynamics and IBM Theory \textbullet\ Oral exam backup slides}
\begin{document}

% ------------------------------------------------------------
\begin{frame}[plain]
\maketutorialtitle
\end{frame}

% ============================================================
% BACKUP SLIDES: Meaning of force spreading + discretization + NS LHS
% ============================================================

% ------------------------------------------------------------
\begin{frame}{[Backup] Force spreading: pointwise meaning}
\small

\textbf{IB coupling definition (force spreading):}
\[
\boxed{
\mathbf{f}(\mathbf{x},t)
=
\int_{\mathcal U}
\mathbf{F}(\mathbf{X},t)\,
\delta\!\big(\mathbf{x}-\boldsymbol{\chi}(\mathbf{X},t)\big)\,d\mathbf{X}
}
\]

\medskip
\textbf{Interpretation in words:}
\begin{itemize}\itemsep3pt
\item $\mathcal U$ is the reference material domain, and $\mathbf{F}(\mathbf{X},t)$ is a force density with respect to its material measure $d\mathbf X$.
\item $\delta(\mathbf{x}-\chi(\mathbf{X},t))$ deposits that force at the \emph{current physical location}
      $\mathbf{x}=\chi(\mathbf{X},t)$.
\item $\mathbf{f}(\mathbf{x},t)$ is the corresponding \emph{Eulerian body-force density} used in Navier--Stokes.
\end{itemize}

\medskip
\textbf{Key idea:} this compactly expresses ``forces carried by the structure become forces applied to the fluid
at the structure's current location.''

\medskip
{\scriptsize\color{DarkGray}Notation: these backup slides use the continuum-mechanics convention of Module~4, with $\mathbf X$ the material (reference) label and $\boldsymbol\chi(\mathbf X,t)$ its current position. Decks 1-2 and 1-3 write the same map in Peskin's convention as $\mathbf X(\mathbf q,t)$, with label $\mathbf q$.\par}

\end{frame}

% ------------------------------------------------------------
\begin{frame}{[Backup] Force spreading over a region}
\footnotesize

Let $\Gamma\subset\Omega$ be any Eulerian region.

\medskip
Start with the definition:
\[
\mathbf{f}(\mathbf{x},t)=\int_{\mathcal U}\mathbf{F}(\mathbf{X},t)\,
\delta(\mathbf{x}-\chi(\mathbf{X},t))\,d\mathbf{X}.
\]

\medskip
Integrate over $\Gamma$:
\[
\int_{\Gamma}\mathbf{f}(\mathbf{x},t)\,d\mathbf{x}
=
\int_{\mathcal U}\mathbf{F}(\mathbf{X},t)\,
\left(\int_{\Gamma}\delta(\mathbf{x}-\chi(\mathbf{X},t))\,d\mathbf{x}\right)d\mathbf{X}.
\]

\medskip
Use the defining property of $\delta$:
\[
\int_{\Gamma}\delta(\mathbf{x}-\chi(\mathbf{X},t))\,d\mathbf{x}
=
\begin{cases}
1, & \chi(\mathbf{X},t)\in\Gamma,\\
0, & \chi(\mathbf{X},t)\notin\Gamma.
\end{cases}
\]

\medskip
\textbf{Conclusion (the intuitive identity):}
\[
\boxed{
\int_{\Gamma}\mathbf{f}(\mathbf{x},t)\,d\mathbf{x}
=
\int_{\{\mathbf{X}:\,\chi(\mathbf{X},t)\in\Gamma\}}\mathbf{F}(\mathbf{X},t)\,d\mathbf{X}.
}
\]

\end{frame}

% ------------------------------------------------------------
\begin{frame}{[Backup] Distributional meaning and Fubini}
\small

\textbf{Subtlety:} $\delta$ is a \emph{distribution}, not an $L^1$ function, so classical Fubini/Tonelli
does not literally apply to
\[
\int_{\Gamma}\int_{\mathcal U}\mathbf{F}(\mathbf{X},t)\,\delta(\mathbf{x}-\chi(\mathbf{X},t))\,d\mathbf{X}\,d\mathbf{x}.
\]

\medskip
\textbf{Two standard justifications used in IB:}

\medskip
\textbf{(A) Weak/distribution meaning (clean):}
for any smooth test function $\boldsymbol{\varphi}(\mathbf{x})$,
\[
\int_{\Omega}\mathbf{f}(\mathbf{x},t)\cdot\boldsymbol{\varphi}(\mathbf{x})\,d\mathbf{x}
:=
\int_{\mathcal U}\mathbf{F}(\mathbf{X},t)\cdot\boldsymbol{\varphi}(\chi(\mathbf{X},t))\,d\mathbf{X}.
\]
No ``swap'' is needed; the delta notation is shorthand for this identity.

\medskip
\textbf{(B) Regularized delta (what we compute):}
replace $\delta$ by $\delta_h$ (bounded, integrable, compactly supported). Then the double integral is absolutely integrable,
and Fubini is valid.

\end{frame}

% ------------------------------------------------------------
\begin{frame}{[Backup] Regularized delta kernel $\delta_h$}
\small

In computation we use a regularized kernel $\delta_h$ supported on a few grid cells:
\[
\mathbf{f}_h(\mathbf{x},t)=\int_{\mathcal U}\mathbf{F}(\mathbf{X},t)\,
\delta_h(\mathbf{x}-\chi(\mathbf{X},t))\,d\mathbf{X}.
\]

\medskip
\textbf{Why this helps:}
\begin{itemize}\itemsep3pt
\item $\delta_h$ is a bounded, integrable function rather than a distribution, so $\mathbf{f}_h$ is an ordinary function.
\item The coupling becomes weighted sums over nearby grid points / Lagrangian points.
\item Partition of unity and first-moment conditions support exact transfer of net force and torque when the discrete weights are compatible; other kernel conditions reduce sensitivity to grid position.
\end{itemize}

\medskip
\textbf{If asked “what does $\delta_h$ do?”:}
it spreads (or interpolates) between the Lagrangian mesh and Eulerian grid over a small neighborhood of width $\sim h$.

\end{frame}

% ------------------------------------------------------------
\begin{frame}{[Backup] One explicit IB time step}
\footnotesize
\setlength{\abovedisplayskip}{3pt}\setlength{\belowdisplayskip}{3pt}%

\textbf{Grids:}
Eulerian grid points $\mathbf{x}_{i}$ (Cartesian mesh spacing $h$), and Lagrangian points $\mathbf{X}_k$
(or FE quadrature points), with positions $\chi_k^n\approx \chi(\mathbf{X}_k,t^n)$.

\medskip
\textbf{One illustrative explicit-coupling time step:}
\begin{enumerate}\itemsep2pt
\item \textbf{Structure forces:} compute $\mathbf{F}_k^n$ from the structure model (springs/beam/FE stresses, etc.).
\item \textbf{Spread to fluid:} for each Eulerian grid point $\mathbf{x}_i$,
\[
\mathbf{f}_i^n \approx \sum_k \mathbf{F}_k^n\,\delta_h(\mathbf{x}_i-\chi_k^n)\,w_k,
\]
where $w_k$ is the Lagrangian quadrature weight.
\item \textbf{Fluid solve:} advance Navier--Stokes using $\mathbf{f}^n$ as a body force to get $\mathbf{u}^{n+1},p^{n+1}$.
\item \textbf{Interpolate to structure:} for each Lagrangian point,
\[
\mathbf U_k^{n+1} \approx \sum_i \mathbf{u}_i^{n+1}\,\delta_h(\mathbf{x}_i-\chi_k^{n})\,v_i,
\]
where $v_i=h^d$ on a uniform grid.
\item \textbf{Update structure:} $\chi_k^{n+1}=\chi_k^{n}+\Delta t\,\mathbf U_k^{n+1}$ (or implicit/semi-implicit variant).
\end{enumerate}
\smallskip
With compatible weights, interpolation is chosen as the discrete adjoint of spreading so the two grids exchange equal power.

\end{frame}

% ------------------------------------------------------------
\begin{frame}{[Backup] Density times material acceleration}
\small

For an incompressible Newtonian fluid with constant $\rho$, constant $\mu$, and body force $\mathbf{f}$:
\[
\rho\left(\frac{\partial \mathbf{u}}{\partial t} + (\mathbf{u}\cdot\nabla)\mathbf{u}\right)
=
-\nabla p + \mu\nabla^2\mathbf{u} + \mathbf{f},
\qquad \nabla\cdot\mathbf{u}=0.
\]

\medskip
\textbf{Left-hand side (LHS):}
\[
\boxed{
\rho\left(\frac{\partial \mathbf{u}}{\partial t} + (\mathbf{u}\cdot\nabla)\mathbf{u}\right)
=
\rho\,\frac{D\mathbf{u}}{Dt}
}
\]
is density times the \textbf{material (total) acceleration} of a fluid parcel.

\medskip
\textbf{Two pieces:}
\begin{itemize}\itemsep3pt
\item \textbf{Local acceleration} $\displaystyle \frac{\partial \mathbf{u}}{\partial t}$:
change in velocity at a fixed point in space.
\item \textbf{Advective acceleration} $\displaystyle (\mathbf{u}\cdot\nabla)\mathbf{u}$:
change in velocity because the parcel moves through spatial velocity gradients.
\end{itemize}

\end{frame}

% ------------------------------------------------------------
\begin{frame}{[Backup] Advective acceleration in components}
\footnotesize

Let $\mathbf{u}=(u,v,w)$ in 3D (or omit $w$ in 2D). Then
\[
(\mathbf{u}\cdot\nabla)\mathbf{u}
=
\left(u\frac{\partial}{\partial x}+v\frac{\partial}{\partial y}+w\frac{\partial}{\partial z}\right)
\begin{pmatrix}u\\v\\w\end{pmatrix}.
\]

\medskip
So each component has the form:
\[
\begin{aligned}
\text{x-component:}\quad
&\frac{Du}{Dt}=\frac{\partial u}{\partial t} + u\frac{\partial u}{\partial x}+v\frac{\partial u}{\partial y}+w\frac{\partial u}{\partial z},\\[3pt]
\text{y-component:}\quad
&\frac{Dv}{Dt}=\frac{\partial v}{\partial t} + u\frac{\partial v}{\partial x}+v\frac{\partial v}{\partial y}+w\frac{\partial v}{\partial z},\\[3pt]
\text{z-component:}\quad
&\frac{Dw}{Dt}=\frac{\partial w}{\partial t} + u\frac{\partial w}{\partial x}+v\frac{\partial w}{\partial y}+w\frac{\partial w}{\partial z}.
\end{aligned}
\]

\medskip
\textbf{Quick oral-exam style interpretation:}
``the parcel accelerates either because time changes at a point (local) or because it advects into a region
with different velocity (advective).''

\vfill
{\tiny\color{black!55}Claude (Anthropic) was used to assist with formatting, organization, and clarification of the instructor's original notes.}

\end{frame}

% ------------------------------------------------------------
\begin{frame}{References for the backup slides}
\footnotesize
\begin{itemize}\itemsep5pt
\item C.~S. Peskin, ``The immersed boundary method,'' \emph{Acta Numerica} 11 (2002), 479--517.
  \href{https://doi.org/10.1017/S0962492902000077}{doi:10.1017/S0962492902000077}.
  (Force spreading, velocity interpolation, regularized delta kernels and their moment conditions, the explicit time step.)
\item B.~E. Griffith and X.~Y. Luo, ``Hybrid finite difference/finite element immersed boundary method,''
  \emph{Int. J. Numer. Meth. Biomed. Eng.} 33 (2017), e2888.
  \href{https://doi.org/10.1002/cnm.2888}{doi:10.1002/cnm.2888}.
  (Reference-configuration notation $\boldsymbol\chi(\mathbf X,t)$, FE quadrature points.)
\item B.~E. Griffith and N.~A. Patankar, ``Immersed methods for fluid--structure interaction,''
  \emph{Annu. Rev. Fluid Mech.} 52 (2020), 421--448.
  \href{https://doi.org/10.1146/annurev-fluid-010719-060228}{doi:10.1146/annurev-fluid-010719-060228}.
\item D.~J. Acheson, \emph{Elementary Fluid Dynamics}, Oxford University Press (1990), Ch.~1 and 6.
  (Material derivative and the Navier--Stokes left-hand side.)
\end{itemize}
\medskip
{\scriptsize\color{DarkGray}No third-party figures are used in these backup slides. Adapted from the instructor's Module 1 decks (1-2, 1-3, 1-4), which follow the sources above.\par}
\end{frame}
\end{document}
